% ===== 原书 p.2 =====
\section{镜像对称}

镜像对称的源流，可以上溯到对环面上弦理论的研究过程中所发现的这样一个认识：弦是以与点粒子全然不同的方式观察这个世界的。对于在半径为 $2\pi R$ 的圆周 $T=\R/2\pi R\Z$ 上运动的自由粒子，不含时的 Schr\"odinger 方程
\begin{equation}
-\frac{\partial^2}{\partial x^2}\Psi(x)=E\Psi(x)
\tag{1}
\end{equation}
\noindent 的解可借助动量 $m\in\Z$ 写成 $\Psi(x)=\exp\left(\sqrt{-1}mx/R\right)$，此时的能量由 $E=m^2/R^2$ 给出。另一方面，为了描述 $T$ 上经典弦的运动，把能量取为泛函
\begin{equation}
\Map(S^1,T)\to\R,\qquad
\gamma\mapsto\int_{S^1}\left(\frac{d\gamma}{d\theta}\right)^2\frac{d\theta}{2\pi}
\tag{2}
\end{equation}
\noindent 则在与给定的卷数 $w\in\Z\cong\pi_1(T)$ 相应的位形中，能量最小者为 $\gamma(\theta)=Rw\theta$，此时的能量由 $R^2w^2$ 给出。在 $T$ 上运动的量子弦同时具有动量与卷数，其能量为
\begin{equation}
E=\frac{m^2}{R^2}+R^2w^2
\tag{3}
\end{equation}
\noindent 这里，若把半径 $R$ 换成 $1/R$，同时把动量与卷数互换，弦的能量便保持不变。这一论证本质上是非相对论的，而且描述弦的状态时除动量与卷数之外还需要无限多个指定弦振动模式的自然数，在这一点上也不够精确；但即使作更精细的考察，也已知道半径为 $R$ 的圆周上的弦理论与半径为 $1/R$ 的圆周上的弦理论无法区分。这称为 $T$ 对偶（T-duality），
\begin{itemize}[leftmargin=2.2em,itemsep=2pt,topsep=3pt]
\item 是点粒子理论中所没有的、弦理论特有的现象；
\item 是经典理论中所没有的量子现象。
\end{itemize}

更一般地，用 $\R^n$ 的秩 $n$ 离散子群 $N\subset\R^n$ 作商所得的环面记为 $T:=\R^n/N$，则在 $T$ 上运动的弦，除了取值于 $T$ 的特征标群 $M=\Hom(T,S^1)\cong\Hom(N,\Z)$ 的动量之外，还具有取值于 $N\cong\pi_1(T)$（即 $M$ 的对偶群）的卷数；弦的动力学在「互换动量与卷数、同时互换环面 $T=N_{\R}/N$ 与对偶环面 $\check T=M_{\R}/M$」这一操作下不变。换言之，对弦理论而言，环面 $T$ 与对偶环面 $\check T$ 无从区分。

超弦理论只存在于 $10$ 维，分为 I 型超弦理论、IIA 型超弦理论、IIB 型超弦理论、$SO(32)$ 杂化弦理论、$E_8\times E_8$ 杂化弦理论五类。其中 $E_8\times E_8$ 杂化弦理论放在实四维时空 $\R^4$ 与复三维 Calabi--Yau 流形 $M$ 的直积 $\R^4\times M$ 上考虑，即可给出接近粒子物理标准模型的模型，这被发现之后，三维 Calabi--Yau 流形上的弦理论的研究便兴盛起来。这里，Calabi--Yau 流形指 Ricci 曲率为零的 K\"ahler 流形。这与全纯体积形式 $\Omega$ 和 K\"ahler 形式 $\omega$ 之间成立关系
\begin{equation}
\frac{1}{n!}\omega^n=(-1)^{n(n-1)/2}\left(\frac{\sqrt{-1}}{2}\right)^n\Omega\wedge\overline{\Omega}
\tag{4}
\end{equation}
% [Tutor Note: (4) 不是 Calabi--Yau 的定义，而是「Ricci 曲率为零」这一条件的体积形式写法：$n$ 维 K\"ahler 流形 Ricci 平坦当且仅当某个全纯 $(n,0)$ 形式的模长平方为常数，调整常数即得 (4)]

% ===== 原书 p.3 =====
等价。粗略地说，镜像对称是指：对任意 Calabi--Yau 流形 $Y$，存在另一个称为其镜像流形的 Calabi--Yau 流形 $\check Y$，使得 $Y$ 上的 IIA 型超弦理论与 $\check Y$ 上的 IIB 型超弦理论在适当意义下同构。「$Y$ 上的 IIA 型超弦理论」与「$\check Y$ 上的 IIB 型超弦理论」这些概念在数学上尚未被严格地表述，因此要把镜像对称表述为数学猜想，就必须使用这样的对象：它在数学上可以被严格表述，并且被期待提取出了超弦理论所含信息的一部分。其典型例子便是所谓的拓扑弦论，而其中的一部分又被以 Gromov--Witten 不变量、量子上同调理论的形式严格地数学化。在这一脉络下关于镜像对称的猜想中最著名的，是经典镜像对称（classical mirror symmetry）：它断言对任意 Calabi--Yau 流形 $Y$ 存在某个 Calabi--Yau 流形 $\check Y$，使得 $Y$ 的 Gromov--Witten 不变量与 $\check Y$ 的周期积分（Hodge 结构的形变）之间有奇妙的关系。这一猜想已由 Givental 以及 Lian--Liu--Yau 等人对环面流形中作为完全交（complete intersection）给出的 Calabi--Yau 流形加以证明，但其证明基于关于环面作用的等变上同调的局部化，这种关系之所以存在的根本理由仍是个谜。

为给镜像对称提供概念性的理解，Kontsevich [Kon95] 所提出的，是同调镜像对称（homological mirror symmetry）。它断言对任意 Calabi--Yau 流形 $Y$ 存在某个 Calabi--Yau 流形 $\check Y$，使得 $Y$ 的 Fukaya 范畴的导出范畴与 $\check Y$ 的相干层的导出范畴作为三角范畴等价：
\begin{equation}
D^b\mathrm{Fuk}\,Y\cong D^b\mathrm{coh}\,\check Y.
\tag{5}
\end{equation}
% [Tutor Note: 原文把 $D^b$ 与 Fuk、coh 分开排印（$D^b$Fuk $Y$），义为 Fuk、coh 分别是 $\mathrm{Fuk}$、$\mathrm{coh}$ 的正体字记号；这里按通常写法统一为 $D^b\mathrm{Fuk}$、$D^b\mathrm{coh}$]
由同调镜像对称推出经典镜像对称，目前仍是猜想。

同调镜像对称提供了把三角范畴看作空间并研究其几何这一崭新的视角。另一方面，作为对镜像对称更具几何性的进路，Strominger--Yau--Zaslow [SYZ96] 提出了如下图景：任意 Calabi--Yau 流形都具有特殊 Lagrange 环面纤维化，把该纤维化中作为纤维的环面换成对偶环面，便得到镜像流形。这常被称为 SYZ 猜想，其精确表述与证明是 Calabi--Yau 流形几何中的中心问题之一。这里，Calabi--Yau 流形的子流形 $L$ 称为特殊 Lagrange 子流形（special Lagrangian submanifold），是指 $L$ 是 Lagrange 子流形，且全纯体积形式的虚部限制到 $L$ 上恒等于零；即 $\Im\Omega|_L=0$。全纯体积形式的实部是校准形式（calibration），特殊 Lagrange 子流形关于这个校准形式是校准子流形（calibrated submanifold）。特别地，特殊 Lagrange 子流形是极小子流形。对 K3 曲面而言，「$L$ 是特殊 Lagrange 子流形」与「$L$ 经超 K\"ahler 旋转后成为复子流形」这两种说法等价。在这个意义上，特殊 Lagrange 环面纤维化是椭圆曲面概念向实的方向与高维方向的推广，而 SYZ 猜想可以看作对 Calabi--Yau 流形给出某种结构定理的猜想。此外，如后文所述，人们也期待由 SYZ 猜想，经由沿纤维的某种 Fourier 变换（称为 SYZ 变换，与族的 Floer 上同调密切相关）推出同调镜像猜想 (5)。
% [Tutor Note: 原文此处「複素部分多様体」直译为「复部分流形」，指复子流形；「実および高次元への一般化」的「実」指纤维的实维数方向（椭圆曲面的纤维是实二维环面）]

% CHUNK c01 DONE pages p2–p3
